\section{Conclusion}\label{section:conclusion}

This project investigated Physics-Informed Neural Networks (PINNs) as mesh-free solvers for option-pricing partial differential equations, exploiting the well-known fact that the Black--Scholes equation reduces to a heat equation under a standard change of variables and is therefore a natural target for the smooth, differentiable surrogates that PINNs produce. The work proceeded along three connected stages: a forward solver for the Black--Scholes equation, a forward solver for the Heston stochastic-volatility PDE, and an inverse calibration in which the structural Heston parameters $\kappa$, $\theta$, $\xi$, and $\rho$ were inferred jointly with the network weights $\boldsymbol{w}$ from noisy synthetic prices. The unifying theme across the three experiments is that a single differentiable price surrogate provides both prices and Greeks ``for free'' through the same automatic-differentiation machinery used to train it \parencite{baydin2018automatic,karniadakis2021physics}, and that this machinery also admits parameter calibration without external solvers.

The Black--Scholes experiment served as a controlled benchmark with a known closed-form solution. At the at-the-money evaluation point $S=K=100$ and $\tau=0.5$, the trained PINN predicted $V_{\mathrm{PINN}}=6.927159$ against the analytical value $V_{\mathrm{BS}}=6.888729$, giving an absolute error of $0.038430$ (approximately $0.56\%$ relative) and a local PDE residual of $2.746582\times 10^{-3}$ after $50{,}000$ Adam epochs (12\,min 8\,s on a single H100). The automatically differentiated Greeks (Delta, Theta, Gamma, Charm, Speed, Color) were qualitatively consistent with the analytical behavior of an at-the-money call. For a problem with a closed-form solution, the value of the PINN approach is methodological rather than computational: the closed form is faster and more accurate, but the same machinery extends seamlessly to models without closed forms or to settings where Greeks across a full surface are needed without finite-difference approximations.

The Heston experiment extended the same architecture to a three-dimensional problem $(S, v, \tau)$ governed by a PDE containing a mixed second derivative $\rho \xi v S \partial^2 V/\partial S\partial v$ and a variance-dynamics drift. Boundary conditions were imposed on both $S=0$ and $S=S_{\max}$ (with $V(S_{\max}, v, \tau)\approx S_{\max}-Ke^{-r_r\tau}$), and the network was trained for $50{,}000$ epochs (24\,min 40\,s on an A100). On the variance slice $v=v_0=0.04$ at $S=K=100$, $\tau=0.5$, the PINN predicted $V_{\mathrm{PINN}}=6.806941$ against the semi-analytical Heston value $V_{\mathrm{Heston}}=6.837137$, an absolute error of $0.030196$ (approximately $0.44\%$). However, the local PDE residual was substantially larger in magnitude ($-1.790738\times 10^{-1}$) and the loss history was visibly less stable than in the Black--Scholes case, consistent with the by-now standard observation that PINN optimization landscapes degrade as the PDE becomes stiffer, more multi-scale, or more strongly coupled \parencite{wang2021understanding,wang2022expert,krishnapriyan2021characterizing}.

The inverse experiment recast the Heston PINN as a calibration engine: the four structural parameters $(\kappa, \theta, \xi, \rho)$ were promoted to trainable variables and optimized jointly with $\boldsymbol{w}$ against $200$ noisy synthetic prices, using a two-stage schedule of $25{,}000$ Adam epochs followed by $5{,}000$ L-BFGS iterations (30{,}000 total, 16\,min 17\,s) under loss weights $\lambda_{\mathrm{PDE}}=20$, $\lambda_{\mathrm{IC}}=10$, $\lambda_{\mathrm{BC}}=5$, $\lambda_{\mathrm{data}}=1000$, and a Feller-penalty weight $\lambda_{\mathrm{Feller}}=0.1$. Judged purely on its ability to reproduce observed prices, the model performed extremely well: the final mean absolute pricing error on the $200$ synthetic observations was $0.028566$, the maximum absolute error was $0.218854$, the median relative error was approximately $0.02\%$, and the parity-plot diagnostic gave $R^2 = 0.999999$. By any pricing yardstick the calibration succeeded.

By any \emph{parameter-recovery} yardstick, however, the same calibration failed. The final recovered parameter vector was $(\kappa, \theta, \xi, \rho) = (1.588,\, 0.0266,\, 0.104,\, 0.994)$, against the true data-generating values $(2.0,\, 0.04,\, 0.3,\, -0.7)$. The most striking discrepancy is the correlation: $\rho$ converged to a strongly positive value despite the data being generated with $\rho=-0.7$, the standard equity-market sign. Mechanistically this is not surprising. The parameter $\rho$ enters the Heston PDE only through the cross term $\rho\xi v S \partial^2 V/\partial S\partial v$, and therefore appears in the data only through the implied-volatility skew that this cross term produces; with $200$ noisy points sparsely scattered across a wide $(S,v,\tau)$ region, the skew signal is far too weak to identify $\rho$ from prices alone. Moreover, the product $\xi\rho$ is more identifiable than $\rho$ individually, and even that product is recovered with the wrong sign here: $\xi\rho \approx (0.104)(0.994)\approx 0.10$ recovered against $(0.3)(-0.7)\approx -0.21$ true. This is not a simple sign confusion in $\rho$ alone but rather points to genuine multi-modality in the data-likelihood landscape, in which the optimizer found a price-consistent representation of the observations through a different combination of $\xi$ and $\rho$.

The financial implications of this two-sided result are concrete. First, forward PINNs are already useful: once trained, the network is a fast, differentiable pricing surrogate that delivers prices and a full surface of Greeks through automatic differentiation, which is directly applicable to pricing and hedging workflows. Second, inverse PINN calibration is a promising direction but is not yet competitive with classical calibration as a black-box parameter-recovery tool. Third, the data deficit identified here is much less severe in practice: real-world calibration uses listed quotes across many strikes and maturities---typically dense option chains in implied-volatility space---which is orders of magnitude richer than the $200$ points used here, and it is precisely that density (especially near-the-money skew) that classical calibration exploits to identify $\rho$. Fourth, PINNs are likely to be most valuable in settings where classical approaches are not available, in particular rough-volatility models \parencite{gatheral2018volatility,bayer2016rough,euch2019roughheston}, for which no easily-invertible characteristic function exists; that extension is left to future work.

Several concrete extensions follow directly from the experiments above. (i) Richer collocation and observation designs spanning a full strike--maturity surface rather than $200$ scattered prices, ideally with controlled noise levels and a near-the-money refinement to expose skew. (ii) Bayesian or regularized priors on $\rho$ and $\xi$ that encode the strong economic prior $\rho<0$ for equity underlyings, turning the under-identified calibration problem into a well-posed MAP problem. (iii) Loss functions formulated in implied-volatility space rather than price space, so that the skew signal that identifies $\rho$ is amplified relative to the level signal that dominates a price-MSE loss. (iv) Extension of the forward solver to American and exotic options, where the PINN's mesh-free property pays off most clearly because no rectangular grid exists. (v) Extension to rough Heston \parencite{euch2019roughheston} as the natural target where no closed-form characteristic function is easily inverted and where a differentiable surrogate is structurally hardest to beat.

The overall picture is that PINNs are a powerful and flexible differentiable framework for financial PDEs. Forward pricing is in good shape across both Black--Scholes and Heston, and the same machinery delivers Greeks essentially for free. Inverse calibration is genuinely harder and more delicate: a near-perfect price fit does not imply structural parameter recovery, and this project demonstrates both sides of that trade-off in a single end-to-end pipeline.
